\subsection{Trabajo futuro}

Una primera línea de continuación consiste en extender la estrategia propuesta a las variantes
paralelas de TMFQSfullstate basadas en OpenMP y MPI. Esta extensión permitiría estudiar la
interacción entre compresión, particionamiento del vector de estado, sincronización entre
hilos o procesos y movimiento de datos entre memorias locales y distribuidas, con el fin de
determinar si el ahorro observado en la versión actual se mantiene en escenarios de mayor
escala.

Una segunda línea corresponde al desarrollo de un planificador de parámetros que seleccione de
forma automática la configuración de compresión más adecuada para cada carga de trabajo. En
el estado actual, parámetros como el tamaño de bloque, el número de entradas de caché, el
nivel de compresión, el uso de filtros previos y el número de hilos se fijan antes de la
ejecución. Un mecanismo adaptativo podría ajustar estas decisiones según el circuito, el
tamaño del estado, la tasa de aciertos en caché o la compresibilidad observada durante la
simulación.

Una tercera dirección consiste en incorporar criterios de activación selectiva. Los resultados
de QFT con alta entropía muestran que comprimir todo el vector de estado no siempre es
beneficioso. Por ello, resulta pertinente estudiar estrategias híbridas en las que ciertos
bloques permanezcan densos, otros se compriman y la compresión pueda desactivarse
temporalmente cuando la redundancia del estado no compense el costo de compresión y
descompresión.

También es relevante ampliar la campaña experimental a otras familias de algoritmos y estados
de entrada. Algoritmos variacionales, circuitos aleatorios, simulaciones de sistemas físicos
con distinta localidad y cargas de trabajo en las que la compresibilidad cambie a lo largo de
la ejecución permitirían caracterizar con mayor precisión los regímenes en los que la
estrategia propuesta resulta favorable o desfavorable.

Finalmente, este trabajo puede extenderse mediante la comparación con otras bibliotecas y
políticas de almacenamiento. La evaluación de compresores sin pérdida alternativos, políticas
de reemplazo distintas de LRU y entornos con GPU o memoria distribuida permitiría establecer
si el equilibrio observado con Blosc responde a propiedades generales de la arquitectura
propuesta o a particularidades de la realización concreta estudiada aquí.
